\documentclass[12pt]{article}
\usepackage{amsmath,amssymb}
\usepackage{booktabs}
\usepackage{geometry}
\geometry{margin=1in}
\usepackage{natbib}

\newcommand{\GntR}{\text{GntR}}
\newcommand{\Venus}{\text{Venus}}
\newcommand{\kB}{k_BT}

\begin{document}

\section*{Supplementary Material}
\setcounter{equation}{0}
\renewcommand{\theequation}{S\arabic{equation}}
\setcounter{figure}{0}
\renewcommand{\thefigure}{S\arabic{figure}}

\subsection*{S1. Derivation of transcription and translation kinetic rate expressions}

\subsubsection*{S1.1. Elementary reaction mechanism for transcription}
We model transcription as a four-step process based on the classical work of McClure \cite{mcclure_1980} and as developed previously \cite{adhikari2020effective}. RNA polymerase (RNAP) binds the promoter of gene $j$ to form a closed complex, which isomerizes to an open complex, from which either abortive initiation or productive elongation occurs:
\begin{align}
\mathcal{G}_j + R_X &\rightleftharpoons (\mathcal{G}_j : R_X)_C \label{eq:step1} \\
(\mathcal{G}_j : R_X)_C &\rightarrow (\mathcal{G}_j : R_X)_O \label{eq:step2} \\
(\mathcal{G}_j : R_X)_O &\rightarrow R_X + \mathcal{G}_j \label{eq:step3} \\
(\mathcal{G}_j : R_X)_O &\rightarrow m_j + R_X + \mathcal{G}_j \label{eq:step4}
\end{align}
where $\mathcal{G}_j$ denotes the concentration of gene $j$, $R_X$ is the concentration of \textit{free} RNAP, $(\mathcal{G}_j : R_X)_C$ and $(\mathcal{G}_j : R_X)_O$ are the closed and open complex concentrations, and $m_j$ is the mRNA product. The key idea is that RNAP acts as a pseudo-enzyme: it binds its substrate (the gene), processes it, and dissociates, analogous to Michaelis-Menten kinetics.

\subsubsection*{S1.2. Steady-state approximation for complexes}
Let $k_+$ (conc$^{-1}$ t$^{-1}$) and $k_-$ (t$^{-1}$) denote the on/off rate constants for RNAP binding at the promoter, $k_I$ (t$^{-1}$) the rate constant for open complex formation, $k_A$ (t$^{-1}$) the abortive initiation rate constant, and $k_{E,j}^X$ (t$^{-1}$) the elongation rate constant for gene $j$. The material balances for the closed and open complexes are:
\begin{align}
\frac{d}{dt}(\mathcal{G}_j : R_X)_C &= k_+\,\mathcal{G}_j\,R_X - k_-\,(\mathcal{G}_j : R_X)_C - k_I\,(\mathcal{G}_j : R_X)_C \label{eq:dCC} \\
\frac{d}{dt}(\mathcal{G}_j : R_X)_O &= k_I\,(\mathcal{G}_j : R_X)_C - k_A\,(\mathcal{G}_j : R_X)_O - k_{E,j}^X\,(\mathcal{G}_j : R_X)_O \label{eq:dOC}
\end{align}
Invoking the pseudo-steady-state approximation (justified because complex formation and dissociation are fast compared to mRNA accumulation), we set Eqns.~\ref{eq:dCC}--\ref{eq:dOC} to zero and solve:
\begin{align}
(\mathcal{G}_j : R_X)_C &= \left(\frac{k_+}{k_- + k_I}\right)\mathcal{G}_j\,R_X \equiv K_{X,j}^{-1}\,\mathcal{G}_j\,R_X \label{eq:CC_ss} \\
(\mathcal{G}_j : R_X)_O &= \left(\frac{k_I}{k_A + k_{E,j}^X}\right)(\mathcal{G}_j : R_X)_C \equiv \tau_{X,j}^{-1}\,(\mathcal{G}_j : R_X)_C \label{eq:OC_ss}
\end{align}
where we define the saturation constant $K_{X,j}$ ($\mu$M) and the dimensionless time constant $\tau_{X,j}$:
\begin{equation}
K_{X,j}^{-1} \equiv \frac{k_+}{k_- + k_I}, \qquad \tau_{X,j}^{-1} \equiv \frac{k_I}{k_A + k_{E,j}^X}
\label{eq:K_tau_def}
\end{equation}
The ratio $K_{X,j}^{-1}$ has units of inverse concentration and governs the affinity of RNAP for the promoter. The ratio $\tau_{X,j}^{-1}$ compares the rate of open complex formation ($k_I$) to the combined rates of elongation and abortive initiation ($k_A + k_{E,j}^X$), and controls whether the rate-limiting step is initiation or elongation.

Combining Eqns.~\ref{eq:CC_ss}--\ref{eq:OC_ss}, the open complex concentration is:
\begin{equation}
(\mathcal{G}_j : R_X)_O = K_{X,j}^{-1}\,\tau_{X,j}^{-1}\,\mathcal{G}_j\,R_X
\label{eq:OC_combined}
\end{equation}

\subsubsection*{S1.3. Transcription rate expression}
The rate of mRNA production for gene $j$ is proportional to the open complex concentration:
\begin{equation}
\boxed{r_{X,j} = k_{E,j}^X\,(\mathcal{G}_j : R_X)_O = k_{E,j}^X\,K_{X,j}^{-1}\,\tau_{X,j}^{-1}\,\mathcal{G}_j\,R_X}
\label{eq:rX_general}
\end{equation}
where $R_X$ is the free (unbound) RNAP concentration. This expression states that the transcription rate is the product of the elongation rate constant, the probability of finding an open complex (captured by $K^{-1}\tau^{-1}\mathcal{G}_j$), and the available RNAP.

\subsubsection*{S1.4. Elongation rate constant}
The elongation rate constant $k_{E,j}^X$ is related to the polymerase elongation speed $\dot{v}_X$ (nt/s) and the gene coding length $l_{G,j}$ (nt):
\begin{equation}
k_{E,j}^X = \frac{\dot{v}_X}{l_{G,j}}
\label{eq:kE_def}
\end{equation}
Similarly, the time constant $\tau_{X,j}$ is parameterized as:
\begin{equation}
\tau_{X,j} = \frac{k_{E,j}^X}{k_I}\cdot\hat{\tau}_{X,j}
\label{eq:tau_parameterized}
\end{equation}
where $k_I = 1/t_{\text{init},X}$ is the initiation rate constant (estimated from McClure assays \cite{mcclure_1980}) and $\hat{\tau}_{X,j}$ is a dimensionless gene-specific modifier that is estimated from data.

\subsubsection*{S1.5. Translation rate expression}
An identical derivation applies to translation, with ribosome ($R_L$) replacing RNAP and mRNA ($m_j$) replacing gene ($\mathcal{G}_j$):
\begin{equation}
\boxed{r_{L,j} = k_{E,j}^L\,K_{L,j}^{-1}\,\tau_{L,j}^{-1}\,m_j\,R_L}
\label{eq:rL_general}
\end{equation}
where $k_{E,j}^L = \dot{v}_L / l_{P,j}$ is the translation elongation rate constant, $l_{P,j}$ is the protein length in amino acids, $K_{L,j}$ is the translation saturation constant, and $\tau_{L,j} = (k_{E,j}^L / k_{I,L})\,\hat{\tau}_{L,j}$ is the translation time constant.

\subsubsection*{S1.6. Limiting behavior of the time constant}
The time constant $\tau_{X,j}$ governs whether gene $j$ is elongation-limited or initiation-limited. Assuming abortive initiation is slow ($k_A \ll k_I, k_{E,j}^X$), then $\tau_{X,j} \simeq k_{E,j}^X / k_I$. In the initiation-limited regime ($\tau_{X,j} \gg 1$, i.e., elongation is fast compared to initiation):
\begin{equation}
r_{X,j} \simeq \frac{k_{E,j}^X}{\tau_{X,j}}\,R_{X,T}\left(\frac{\mathcal{G}_j}{K_{X,j} + \mathcal{G}_j}\right)
\label{eq:rX_init_limited}
\end{equation}
which has the familiar Michaelis-Menten form. In the elongation-limited regime ($\tau_{X,j} \ll 1$):
\begin{equation}
r_{X,j} \simeq k_{E,j}^X\,R_{X,T}\left(\frac{\mathcal{G}_j}{K_{X,j}\tau_{X,j} + \mathcal{G}_j}\right)
\label{eq:rX_elong_limited}
\end{equation}
where the effective saturation constant is reduced by $\tau_{X,j}$, reflecting tighter RNAP binding when elongation is slow.

%% ============================================================
\subsection*{S2. Derivation of free machinery concentrations}

\subsubsection*{S2.1. RNAP balance for multiple competing genes}
Consider $\mathcal{N}$ genes competing for a shared pool of RNAP with total concentration $R_{X,T}$. At any time, RNAP is either free or sequestered in closed or open complexes with one of the $\mathcal{N}$ genes:
\begin{equation}
R_{X,T} = R_X + \sum_{j=1}^{\mathcal{N}} \left[(\mathcal{G}_j : R_X)_C + (\mathcal{G}_j : R_X)_O\right]
\label{eq:RNAP_balance}
\end{equation}
Substituting the steady-state expressions (Eqns.~\ref{eq:CC_ss}--\ref{eq:OC_combined}):
\begin{equation}
R_{X,T} = R_X + \sum_{j=1}^{\mathcal{N}} K_{X,j}^{-1}\,\mathcal{G}_j\,R_X\,(1 + \tau_{X,j}^{-1})
= R_X\left[1 + \sum_{j=1}^{\mathcal{N}} \frac{(1 + \tau_{X,j}^{-1})\,\mathcal{G}_j}{K_{X,j}}\right]
\label{eq:RNAP_balance_expanded}
\end{equation}
Solving for the free RNAP concentration:
\begin{equation}
\boxed{R_X = \frac{R_{X,T}}{1 + \displaystyle\sum_{j=1}^{\mathcal{N}} \frac{(\tau_{X,j} + 1)\,\mathcal{G}_j}{\tau_{X,j}\,K_{X,j}}}}
\label{eq:RX_free}
\end{equation}
This expression shows that free RNAP decreases as more genes compete for the shared pool. Each gene $j$ sequesters RNAP in proportion to its concentration $\mathcal{G}_j$, its affinity $K_{X,j}^{-1}$, and the factor $(\tau_{X,j}+1)/\tau_{X,j}$ which accounts for RNAP trapped in both closed and open complexes.

\subsubsection*{S2.2. Free ribosome concentration}
By identical reasoning, the free ribosome concentration is:
\begin{equation}
\boxed{R_L = \frac{R_{L,T}}{1 + \displaystyle\sum_{j=1}^{\mathcal{N}} \frac{(\tau_{L,j} + 1)\,m_j}{\tau_{L,j}\,K_{L,j}}}}
\label{eq:RL_free}
\end{equation}
where $m_j$ is the mRNA concentration of gene $j$, which varies dynamically. Thus $R_X$ is constant in time (since gene concentrations are fixed for linear DNA), while $R_L$ changes as mRNA levels evolve.

\subsubsection*{S2.3. Single-gene rate expression}
For a single gene, substituting Eqn.~\ref{eq:RX_free} into Eqn.~\ref{eq:rX_general} and simplifying gives the closed-form transcription rate:
\begin{equation}
r_{X,j} = k_{E,j}^X\,R_{X,T}\left(\frac{\mathcal{G}_j}{\tau_{X,j}\,K_{X,j} + (1 + \tau_{X,j})\,\mathcal{G}_j}\right)
\label{eq:rX_single}
\end{equation}
To verify: from Eqn.~\ref{eq:rX_general}, $r_{X,j} = k_{E,j}^X\,\mathcal{G}_j\,R_X / (\tau_{X,j}\,K_{X,j})$. For a single gene, Eqn.~\ref{eq:RX_free} gives $R_X = R_{X,T}\,\tau_{X,j}\,K_{X,j} / [\tau_{X,j}\,K_{X,j} + (\tau_{X,j}+1)\,\mathcal{G}_j]$. Substituting and canceling $\tau_{X,j}\,K_{X,j}$ yields Eqn.~\ref{eq:rX_single}. Defining the maximum rate $V_{X,j}^{\max} \equiv R_{X,T}\,({\dot{v}_X}/{l_{G,j}}) = k_{E,j}^X\,R_{X,T}$, this can be written as:
\begin{equation}
r_{X,j} = V_{X,j}^{\max}\left(\frac{\mathcal{G}_j}{\tau_{X,j}\,K_{X,j} + (1 + \tau_{X,j})\,\mathcal{G}_j}\right)
\label{eq:rX_Vmax}
\end{equation}
which is the form presented in the main text.

\subsubsection*{S2.4. Multi-gene competitive allocation}
For multiple genes competing for the same RNAP pool, we define the allocation fraction:
\begin{equation}
f_{X,j} \equiv \frac{\mathcal{G}_j}{\tau_{X,j}\,K_{X,j} + (1 + \tau_{X,j})\,\mathcal{G}_j}
\label{eq:fX_def}
\end{equation}
and the free RNAP as $R_{X,\text{free}} = R_{X,T}/(1 + \sum_j f_{X,j})$. The transcription rate for gene $j$ in the multi-gene system is then:
\begin{equation}
r_{X,j} = R_{X,\text{free}} \cdot k_{E,j}^X \cdot f_{X,j}
\label{eq:rX_with_f}
\end{equation}
or equivalently, $r_{X,j} = V_{X,j}^{\max}\,f_{X,j}$ where now $V_{X,j}^{\max} = R_{X,\text{free}}\,(\dot{v}_X / l_{G,j})$ uses the \textit{free} RNAP concentration rather than the total. This formulation treats RNAP as a shared enzyme that is competitively allocated among genes, analogous to competitive substrate inhibition in enzyme kinetics. The factored form separates the global competition effect (through $R_{X,\text{free}}$, which decreases as more genes sequester RNAP) from the individual gene's kinetics (through $f_{X,j}$, which depends only on gene $j$'s concentration and kinetic parameters).

%% ============================================================
\subsection*{S3. Statistical thermodynamic derivation of the transcription control function}

\subsubsection*{S3.1. Promoter microstates and the partition function}
We model transcriptional regulation using the statistical thermodynamic framework developed by Ackers et al.\ \cite{ackers1982quantitative} and adapted for synthetic circuits by Moon et al.\ \cite{moon2012genetic}. The central idea is that the promoter region of gene $j$ can exist in a finite number of discrete microscopic configurations (microstates), each corresponding to a different combination of bound transcription factors and RNA polymerase. At thermal equilibrium, the probability of finding the promoter in microstate $s$ is given by the Boltzmann distribution:
\begin{equation}
p_s = \frac{W_s\,g_s(\mathbf{P})}{\displaystyle\sum_{s' \in \mathcal{C}_j} W_{s'}\,g_{s'}(\mathbf{P})}
\label{eq:p_microstate}
\end{equation}
where the sum is over all possible configurations $\mathcal{C}_j$ of promoter $j$, and:
\begin{itemize}
\item $W_s = \exp(-\epsilon_s / \kB)$ is the Boltzmann weight (dimensionless) for configuration $s$, determined by the pseudo-energy $\epsilon_s$ of that configuration relative to a reference state (the empty promoter, $\epsilon_0 = 0$, $W_0 = 1$).
\item $g_s(\mathbf{P})$ is a binding function that depends on the concentrations of the relevant transcriptional regulators $\mathbf{P}$. For a configuration in which regulator $i$ is bound, $g_s$ includes a Hill-type occupancy term; for the empty or RNAP-only states, $g_s = 1$.
\end{itemize}

\subsubsection*{S3.2. Control function as probability of productive configurations}
The transcription control function $\bar{u}_j$ is defined as the probability that the promoter is in a configuration that leads to productive transcription (i.e., RNAP is bound and positioned to initiate):
\begin{equation}
\boxed{\bar{u}_j = \frac{\displaystyle\sum_{s \in \chi_j} W_s\,g_s(\mathbf{P})}{\displaystyle\sum_{s \in \mathcal{C}_j} W_s\,g_s(\mathbf{P})}}
\label{eq:u_general}
\end{equation}
where $\chi_j \subseteq \mathcal{C}_j$ is the subset of configurations that are transcriptionally active (i.e., lead to expression). The value $\bar{u}_j \in [0, 1]$ represents the fraction of time the promoter spends in a productive state, and directly modulates the transcription rate: $\dot{m}_j \propto r_{X,j}\,\bar{u}_j$.

The pseudo-energies $\epsilon_s$ encode the thermodynamic favorability of each configuration. A strongly negative $\epsilon_s$ (large $W_s$) indicates a highly favorable binding interaction; a positive $\epsilon_s$ indicates an unfavorable configuration. These pseudo-energies are model parameters estimated from data.

\subsubsection*{S3.3. Application to the gluconate sensor circuit}
The circuit contains three genes, each with a distinct promoter architecture:

\paragraph{P70 promoter (GntR gene).} The P70 promoter can exist in three states: (i) empty, (ii) core RNAP bound (without $\sigma_{70}$), or (iii) holoenzyme ($\sigma_{70}$-RNAP) bound. Both RNAP-bound states are productive. The partition function and control function are:
\begin{equation}
\bar{u}_{\GntR} = \frac{W_{\GntR,\text{cRNAP}} + W_{\GntR,\sigma_{70}}\,b_{\GntR,\sigma_{70}}}{1 + W_{\GntR,\text{cRNAP}} + W_{\GntR,\sigma_{70}}\,b_{\GntR,\sigma_{70}}}
\label{eq:u_GntR}
\end{equation}
where $W_{\GntR,\text{cRNAP}} = \exp(-\epsilon_{\GntR,\text{cRNAP}}/\kB)$ and $W_{\GntR,\sigma_{70}} = \exp(-\epsilon_{\GntR,\sigma_{70}}/\kB)$ are Boltzmann weights for core RNAP and holoenzyme binding, respectively, and:
\begin{equation}
b_{\GntR,\sigma_{70}} = \frac{[\sigma_{70}]^{n_{\GntR,\sigma_{70}}}}{K_{\GntR,\sigma_{70}}^{n_{\GntR,\sigma_{70}}} + [\sigma_{70}]^{n_{\GntR,\sigma_{70}}}}
\label{eq:b_GntR_sigma}
\end{equation}
is a Hill function describing the fractional occupancy of $\sigma_{70}$ on the holoenzyme.

\paragraph{Modified P70 promoter (Venus gene).} The mP70 promoter includes a GntR operator between the $-35$ and $-10$ regions. It can exist in four states: (i) empty, (ii) core RNAP bound, (iii) holoenzyme bound, or (iv) GntR repressor bound to the operator. The GntR-bound state is \textit{not} productive. D-gluconate relieves repression by binding GntR and preventing it from occupying the operator:
\begin{equation}
\bar{u}_{\Venus} = \frac{W_{\Venus,\text{cRNAP}} + W_{\Venus,\sigma_{70}}\,b_{\Venus,\sigma_{70}}}{1 + W_{\Venus,\text{cRNAP}} + W_{\Venus,\sigma_{70}}\,b_{\Venus,\sigma_{70}} + W_{\Venus,\GntR}\,b_{\Venus,\GntR}}
\label{eq:u_Venus}
\end{equation}
The GntR repression term in the denominator reduces $\bar{u}_{\Venus}$ when active (unliganded) GntR is present. The effective GntR concentration available for repression depends on gluconate binding:
\begin{equation}
[\GntR]_{\text{eff}} = (1 - f_{\text{bound}})\,[\GntR]_{\text{total}}, \qquad f_{\text{bound}} = \frac{[\text{gluconate}]^{n_{\text{gluc}}}}{[\text{gluconate}]^{n_{\text{gluc}}} + K_{\text{gluc}}^{n_{\text{gluc}}}}
\label{eq:gluconate_binding}
\end{equation}
where $f_{\text{bound}}$ is the fraction of GntR bound to D-gluconate (and therefore unable to repress), $K_{\text{gluc}}$ is the dissociation constant, and $n_{\text{gluc}}$ is the Hill coefficient for cooperative gluconate binding to the GntR dimer. At saturating gluconate ($f_{\text{bound}} \to 1$), $[\GntR]_{\text{eff}} \to 0$ and repression is fully relieved.

\paragraph{Constitutive promoter ($\sigma_{70}$ gene).} Since $\sigma_{70}$ is supplied exogenously as protein (not encoded on a gene in the circuit), its ``gene'' concentration is zero and no mRNA is produced. However, for completeness, the constitutive promoter control function would be:
\begin{equation}
\bar{u}_{\sigma_{70}} = \frac{W_{\sigma_{70},\text{RNAP}}}{1 + W_{\sigma_{70},\text{RNAP}}}
\label{eq:u_sigma70}
\end{equation}
with $W_{\sigma_{70},\text{RNAP}} = \exp(-0.5)$ fixed.

\subsubsection*{S3.4. Connection between microstates and macroscopic control}
The partition function formalism provides a principled bridge between molecular-level binding events and macroscopic gene expression rates. Each pseudo-energy $\epsilon_s$ encodes the net free energy change for a specific promoter configuration relative to the empty state, integrating contributions from protein-DNA contacts, conformational changes, and steric effects. The Hill functions $b(\cdot)$ capture the concentration dependence of regulator occupancy, with the Hill coefficient $n$ reflecting apparent cooperativity (which may arise from true cooperative binding, conformational switching, or other mechanisms). The resulting control function $\bar{u}_j$ is a smooth, bounded function that modulates the kinetic rate $r_{X,j}$ at each time step, coupling the thermodynamic state of the promoter to the dynamic mass balance equations.

%% ============================================================
\subsection*{S4. Consumable resource depletion}

The second layer of the resource model describes the irreversible depletion of consumable resources (nucleotides, energy cofactors, amino acids) that cannot be regenerated in the cell-free system. We define $\epsilon_X(t) \in [0,1]$ and $\epsilon_L(t) \in [0,1]$ as the fractional availability of transcriptional and translational consumable resources, with $\epsilon_X(0) = \epsilon_L(0) = 1$. The depletion dynamics are:
\begin{align}
\frac{d\epsilon_X}{dt} &= -\alpha_X\left(\sum_{j=1}^{\mathcal{N}} l_{G,j}\,r_{X,j}\,\bar{u}_j\right)\epsilon_X \label{eq:epsX} \\
\frac{d\epsilon_L}{dt} &= -\alpha_L\left(\sum_{j=1}^{\mathcal{N}} l_{P,j}\,r_{L,j}\right)\epsilon_L \label{eq:epsL}
\end{align}
where $\alpha_X$ and $\alpha_L$ are consumption rate constants (units: inverse flux), $l_{G,j}$ is the gene length (nt), and $l_{P,j}$ is the protein length (aa). The depletion rates are proportional to the total biosynthetic flux weighted by sequence length, reflecting the stoichiometric requirement that longer genes consume more nucleotides per transcript and longer proteins consume more amino acids per polypeptide. The resource fractions multiply the production terms in the mRNA and protein balance equations, providing a mechanistic link between cumulative biosynthetic load and the progressive decline of cell-free reaction productivity.

%% ============================================================
\subsection*{S5. Complete model equations}

The full model consists of 11 ordinary differential equations. Gene concentrations are constant (linear DNA, no replication):
\begin{equation}
\frac{d\mathcal{G}_j}{dt} = 0 \qquad j = 1, 2, 3
\end{equation}

The mRNA balances couple transcription (modulated by promoter control and resource availability) with first-order degradation:
\begin{equation}
\frac{dm_j}{dt} = r_{X,j}\,\bar{u}_j\,\epsilon_X - \theta_{m,j}\,m_j \qquad j = 1, 2, 3
\label{eq:mRNA_balance}
\end{equation}
where $\theta_{m,j} = \delta_{m,j}\,\ln 2 / t_{1/2,m}$ is the effective degradation rate constant, $\delta_{m,j}$ is a gene-specific modifier, and $t_{1/2,m}$ is the characteristic mRNA half-life.

The protein balances have an analogous structure:
\begin{equation}
\frac{dp_j}{dt} = r_{L,j}\,\epsilon_L - \theta_{p,j}\,p_j \qquad j = 1, 2, 3
\label{eq:protein_balance}
\end{equation}
where $\theta_{p,j} = \delta_{p,j}\,\ln 2 / t_{1/2,p}$.

Together with the resource depletion equations (Eqns.~\ref{eq:epsX}--\ref{eq:epsL}), the algebraic expressions for transcription control (Eqns.~\ref{eq:u_GntR}--\ref{eq:u_sigma70}), machinery allocation (Eqns.~\ref{eq:RX_free}--\ref{eq:RL_free}), and kinetic rates (Eqns.~\ref{eq:rX_general}, \ref{eq:rL_general}), this defines a closed system of 11 ODEs with 25 estimated parameters.

\bibliographystyle{plos2009}
\bibliography{References}

\end{document}
